\chapter{Component Implementation}

\section{Backend Implementation}

The backend implementation is organized by API surface and business workflow.
Controllers define endpoint paths and dependencies. Services enforce business
rules such as merchant readiness, payment idempotency, refund eligibility,
webhook retry policy, RBAC, audit, and merchant portal scoping. Repositories
hold focused persistence and aggregation queries.

Important implementation points include:

\begin{itemize}
  \item internal auth and merchant portal auth are implemented as separate
  session systems;
  \item merchant-scoped read APIs derive scope from authenticated context rather
  than trusting request parameters;
  \item callback processing and reconciliation are implemented as explicit
  service logic instead of ad hoc controller branches;
  \item write actions that matter operationally emit audit rows with masked
  before and after state snapshots.
\end{itemize}

The Merchant Portal feature introduced a dedicated portal account model,
\texttt{merchant\_users}, and a separate session-auth service. This avoids the
anti-pattern of using HMAC API credentials for human dashboard login. Internal
\texttt{ADMIN} users create and reset merchant portal users from the Ops
Dashboard; generated passwords are returned once and only password hashes are
stored.

\section{Frontend Implementation}

The Ops Dashboard and Merchant Dashboard are separate React, Vite, and
TypeScript applications. Both use React Router for navigation and TanStack
Query for server state. The Merchant Dashboard analytics page uses Recharts for
interactive charts, tooltips, and compact visual summaries.

Important frontend implementation details include:

\begin{itemize}
  \item list and detail layouts for operational scanning on desktop;
  \item responsive one-column behavior on smaller screens;
  \item loading, error, and empty states for async data;
  \item scrollable long lists for recent payments, audit logs, portal users,
  webhooks, credentials, and chart data tables;
  \item stable status and session panels so context does not jump between pages;
  \item explicit label associations on the merchant login form for
  accessibility.
\end{itemize}

\section{Asynchronous Service Implementation}

The current implementation includes the key asynchronous business operations as
callable application services:

\begin{itemize}
  \item payment expiration service for overdue pending payments;
  \item webhook event factory for payment and refund final states;
  \item webhook delivery service with retry scheduling and attempt persistence;
  \item reconciliation services for mismatched or late provider evidence;
  \item internal manual retry endpoint for failed webhook events.
\end{itemize}

However, the project does not yet include a standalone long-running scheduler or
worker runtime as part of the deployed stack. The implemented services are
ready to be invoked by a scheduler or worker trigger, and this boundary is
tested in the backend suite.

\section{Demo Data Implementation}

The demo seed script is explicit and idempotent. It creates a demo merchant,
portal user, credential metadata, payments, refund rows, webhooks, delivery
attempts, and callback evidence. The script does not run automatically on
deploy because production-like deployments should not silently create demo
users.

\begin{lstlisting}[language=bash,caption={Manual demo seed command}]
cd backend
python scripts/seed_dashboard_demo.py
\end{lstlisting}

For sandbox container deployment, the equivalent command is:

\begin{lstlisting}[language=bash,caption={Sandbox demo seed command}]
cd /opt/mini-payment-gateway
docker compose -f docker-compose.sandbox.yml run --rm backend \
  python scripts/seed_dashboard_demo.py
\end{lstlisting}
